\section*{第 \ref{ch:b1:phase-changes} 章——纯物质的相变}

\begin{solution}{exo:b1:phase-changes:1}
(a) 蒸气：\qty{0.5}{kPa} 低于 $P_s(\qty{20}{\degreeCelsius}) = \qty{2.3}{kPa}$。
(b) 固体。(c) 高于 $T_c = \qty{374}{\degreeCelsius}$：超临界流体，
不存在液气之分。(d) \omterm{prop:b1:phase-changes:PTdiagram}{三相点}：三相共存。水洼
从表面蒸发，只要空气中的蒸气压
低于 $P_s(\qty{20}{\degreeCelsius})$；沸腾只发生在
液体内部能形成气泡的情形。
\end{solution}

\begin{solution}{exo:b1:phase-changes:2}
$0.1 \times 2257 + 0.1 \times 4.18\,(100 - T_f) = 1.0 \times 4.18\,(T_f - 20)$：
$351.1 = 4.60\,T_f$，$T_f = \qty{76}{\degreeCelsius}$。若改用热水
则 $100 - T_f = 10(T_f - 20)$，$T_f = \qty{27}{\degreeCelsius}$——
\omterm{def:b1:phase-changes:phase}{潜热}相当于 \qty{540}{K} 的显热。
\end{solution}

\begin{solution}{exo:b1:phase-changes:3}
$v = \qty{2.0e-3}{m^3/kg}$；$x = (2.0 - 1.09) \times 10^{-3}/0.392 = 2.3 \times 10^{-3}$：
\qty{2.3}{g} 的蒸气充满 \qty{0.91}{L}；\qty{0.998}{kg} 的液体
占有体积 $0.998 \times 1.09 = \qty{1.09}{L}$。
\end{solution}

\begin{solution}{exo:b1:phase-changes:4}
$2257/373 = \qty{6.05}{kJ/(kg.K)}$，即 $\times 0.018 = \qty{109}{J/(mol.K)}$；
熔化 $334/273 \times 0.018 = \qty{22}{J/(mol.K)}$。水的 $109$ 明显
高于特鲁顿规则的 $88$：氢键使液态水比普通液体更有序，
因此汽化时无序度增加得更多。
\end{solution}

\begin{solution}{exo:b1:phase-changes:5}
$1/T = 1/373.15 - \ln(P/P_0)/4890$。珠峰：$\ln(0.33/1.013) = -1.12$，
$1/T = 2.680 \times 10^{-3} + 2.29 \times 10^{-4}$，$T = \qty{344}{K} =
\qty{71}{\degreeCelsius}$。在 \qty{0.80}{bar} 时：$\ln = -0.236$，$T = \qty{366}{K} =
\qty{93}{\degreeCelsius}$。在 \qty{71}{\degreeCelsius} 时，蛋清凝固得很慢，
蛋黄几乎不凝固：\qty{3}{min} 的鸡蛋变成
\qty{10}{min} 甚至更久的鸡蛋。
\end{solution}

\begin{solution}{exo:b1:phase-changes:6}
$\ln(1.8/1.013) = 0.575$：$1/T = 2.680 \times 10^{-3} - 1.18 \times 10^{-4}$，
$T = \qty{390}{K} = \qty{117}{\degreeCelsius}$；在 \qty{2.5}{bar} 时：$\ln = 0.903$，
$T = \qty{401}{K} = \qty{128}{\degreeCelsius}$。蒸汽：$\qty{2}{L}/0.98 =
\qty{2.0}{g}$——高压锅内几乎只有水和少量
蒸汽。
\end{solution}

\begin{solution}{exo:b1:phase-changes:7}
$\dd P/\dd T = 3.34 \times 10^5/[273 \times (-9.1 \times 10^{-5})] = \qty{-1.3e7}{Pa/K}
= \qty{-130}{bar/K}$：冰在
约 \qty{130}{bar} 的压强下于 \qty{-1}{\degreeCelsius} 熔化。在 \qty{3}{km} 厚的冰层之下，$P = 917 \times 9.8 \times 3000 =
\qty{270}{bar}$：熔点降至 \qty{-2}{\degreeCelsius}；冰川底部可以处于
熔点，而表面远比它冷，并且由
融水形成的薄膜使冰盖得以滑动。
\end{solution}

\begin{solution}{exo:b1:phase-changes:8}
$W = -P(v_v - v_l) = -1.013 \times 10^5 \times 1.672 = \qty{-169}{kJ}$（水把
大气推回）；$Q = L_v = \qty{2257}{kJ}$；$\Delta U = 2257 - 169
= \qty{2088}{kJ}$；$\Delta H = \qty{2257}{kJ}$；$\Delta S = 2257/373 = \qty{6.05}{kJ/K}$。
若\omterm{def:b1:heat-engines:machine}{热源}为 \qty{1500}{K} 的火焰：$S_{\text{交换}} = 2257/1500 = \qty{1.50}{kJ/K}$，
产生熵 \qty{4.5}{kJ/K}——火焰的大部分品质都浪费在加热
水上。
\end{solution}

\begin{solution}{exo:b1:phase-changes:9}
$P_s(\qty{22}{\degreeCelsius}) \approx 2.34 + 0.4 \times 0.83 = \qty{2.67}{kPa}$，
$P_v = 0.70 \times 2.67 = \qty{1.87}{kPa}$；\omterm{ex:b1:phase-changes:dew}{露点}满足 $P_s = 1.87$：$15 +
5 \times 0.17/0.64 = \qty{16}{\degreeCelsius}$。质量：$\rho_v = 1870 \times 0.018/
(8.314 \times 295) = \qty{13.7}{g/m^3}$，房间内共 \qty{0.69}{kg}。玻璃窗温度为
\qty{12}{\degreeCelsius}：$P_s \approx \qty{1.42}{kPa} < 1.87$，窗户起雾。在
\qty{25}{\degreeCelsius} 时，湿度降至 $1.87/3.17 = 59\%$。
\end{solution}

\begin{solution}{exo:b1:phase-changes:10}
(a) $h$ 是状态函数：经固 $\to$ 液 $\to$ 汽（在
\omterm{prop:b1:phase-changes:PTdiagram}{三相点}处）或由固 $\to$ 汽直接转变，得到的 $\Delta h$ 相同：
$L_s = L_f + L_v$。(b) 由于 $v_s, v_l \ll v_v$，两条斜率均为 $L/(Tv_v)$，
在相同的 $T$ 与 $v_v$ 下，二者之比为 $L_s/L_v > 1$。(c) $v_v = RT/(MP) =
8.314 \times 273.16/(0.018 \times 611) = \qty{206}{m^3/kg}$；汽化为
$2.50 \times 10^6/(273.16 \times 206) = \qty{44}{Pa/K}$，升华为 $2.83 \times 10^6/
(273.16 \times 206) = \qty{50}{Pa/K}$。(d) 在 \qty{0}{\degreeCelsius} 以下，
\qty{1}{bar} 下不存在液体；若空气中的蒸气压低于
冰面上的 $P_s$（\qty{-10}{\degreeCelsius} 时为 \qty{260}{Pa}），则雪
升华——干冷的风把雪吹走。
\end{solution}

\begin{solution}{exo:b1:phase-changes:11}
(a) $P_s(\qty{105}{\degreeCelsius}) \approx 1.013 + 5 \times 0.036 = \qty{1.19}{bar}$：
$\Delta P = \qty{0.18}{bar}$，$r^\ast = 2 \times 0.059/(1.8 \times 10^4) = \qty{6.6}{\micro m}$。
(b) $\Delta P = \qty{3.6}{kPa}$，$r^\ast = \qty{33}{\micro m}$。(c) $x = c\,\Delta T/L_v =
4.18 \times 5/2257 = 0.93\%$：每升水约 \qty{9.3}{g}，即 $9.3 \times 10^{-3}
\times 1.67 = \qty{16}{L}$ 的蒸汽，在极短时间内于
杯内产生——水被喷出，即沸腾。(d) 薯片的
孔隙中含有空气：已存在数十微米量级的气泡，因此
只要 $P_s$ 达到 $P_0$，沸腾便开始。
\end{solution}

\begin{solution}{exo:b1:phase-changes:12}
(a) 节流阀之后 $h = \qty{256}{kJ/kg}$，温度为 \qty{-10}{\degreeCelsius}：$x =
(256 - 187)/(392 - 187) = 0.34$。(b) 蒸发器 $q_c = 392 - 256 =
\qty{136}{kJ/kg}$；压缩机 $w = 430 - 392 = \qty{38}{kJ/kg}$；冷凝器
$q_h = 256 - 430 = \qty{-174}{kJ/kg}$；$136 + 38 - 174 = 0$。(c) $e_p =
174/38 = 4.6$，$e_f = 136/38 = 3.6$；卡诺 $313/50 = 6.3$ 和 $263/50 =
5.3$：约为理想值的 $0.7$。(d) $\dot m = 6000/174000 = \qty{34}{g/s}$，
每分钟 \qty{2.1}{kg}。
\end{solution}

\begin{solution}{pb:b1:phase-changes:1}
\textbf{1.} 在 \qty{1}{atm} 下，水于 \qty{100}{\degreeCelsius} 沸腾；只要
仍有液体，所供给的任何热量都用来产生该温度下的蒸气
（\omterm{def:b1:phase-changes:phase}{潜热}），而不是使液体升温：温度
由压强决定，保持不变。

\textbf{2.} $\dd P_s/\dd T = L_v/[T(v_v - v_l)]$，其中 $v_l \ll v_v = RT/(MP_s)$
（\omterm{def:b1:kinetic-theory:model}{理想气体}蒸气）：$\dd P_s/P_s = (ML_v/R)\,\dd T/T^2$；取 $L_v$
为常量，积分得 $\ln(P_s/P_0) = (ML_v/R)(1/T_0 - 1/T)$。

\textbf{3.} $1/T = 1/373.15 - \ln(1.8/1.013)/4890 = 2.562 \times 10^{-3}$：$T =
\qty{390}{K} = \qty{117}{\degreeCelsius}$。

\textbf{4.} $m = \Delta P\,A/g = 0.8 \times 10^5 \times 2.0 \times 10^{-5}/9.8 =
\qty{0.16}{kg}$。

\textbf{5.} $Q = 2.0 \times 4.18 \times 97 = \qty{811}{kJ}$；$t = 811/2.0 =
\qty{405}{s} \approx \qty{7}{min}$。

\textbf{6.} $2000/(2.2 \times 10^6) = \qty{0.91}{g/s} = \qty{55}{g/min}$——两
升在 \qty{37}{min} 内用完。把电炉调到
使阀门仍保持嘶响的最小档位：温度由压强决定，
而不取决于火焰。

\textbf{7.} $+\qty{17}{K}$：快 $2^{1.7} = 3.2$ 倍，\qty{30}{min}
$\to$ \qty{9}{min}。珠峰：$-\qty{30}{K}$，$2^{-3}$，需四小时。

\textbf{8.} \omterm{ex:b1:phase-changes:dew}{相对湿度} $= P_v/P_s(T)$。$P_v = 0.60 \times 3.17 =
\qty{1.90}{kPa}$；$\rho_v = P_vM/(RT) = 1900 \times 0.018/(8.314 \times 298) =
\qty{13.8}{g/m^3}$。

\textbf{9.} $P_s(T_d) = \qty{1.90}{kPa}$：介于 \qty{15}{\degreeCelsius}
（\qty{1.70}{kPa}）与 \qty{20}{\degreeCelsius}（\qty{2.34}{kPa}）之间：$T_d = 15 + 5 \times 0.20/0.64
= \qty{16.6}{\degreeCelsius}$。夜间向夜空辐射散热的草地，或
从冰箱取出的瓶子，温度低于 $T_d$：与之接触的蒸气
则超过饱和而凝结。

\textbf{10.} $TP^{(1-\gamma)/\gamma} = \text{常量}$：$\dd T/T = \frac{\gamma - 1}{\gamma}
\dd P/P$；$\dd P = -\rho g\,\dd z$，其中 $\rho = PM/(RT)$：$\dd T/T = -\frac{\gamma -
1}{\gamma}\frac{Mg}{RT}\dd z$，故 $\dd T/\dd z = -\frac{\gamma - 1}{\gamma}\frac{Mg}{R} =
-g/c_p$，由于 $c_p = \gamma R/[(\gamma - 1)M]$。取值 $9.8/1000 = \qty{9.8}{K/km}$。

\textbf{11.} $T - T_d$ 以 $9.8 - 2 = \qty{7.8}{K/km}$ 的速率减小：$z = (25 - 16.6)/
7.8 = \qty{1.1}{km}$。

\textbf{12.} $1 \times 2.45 = \qty{2.45}{kJ}$（每千克空气）：$\Delta T =
\qty{2.45}{K}$。释放的热量抵消了部分绝热降温：一团饱和气块
仅以 \qtyrange{5}{6}{K/km} 降温，比周围空气更暖、
更轻，并继续上升——\omterm{def:b1:phase-changes:phase}{潜热}正是
雷暴的燃料。

\textbf{13.} $10^9 \times \qty{1}{g} = \qty{1000}{t}$；$10^6 \times 2.45 \times 10^6 =
\qty{2.5e12}{J} = \qty{0.68}{GW h}$——相当于一座大型发电厂四十分钟的出力。

\textbf{14.} $P_v = 0.80 \times 3.17 = \qty{2.54}{kPa}$；$P_s(\qty{18}{\degreeCelsius})
\approx 1.70 + 0.6 \times 0.64 = \qty{2.08}{kPa} < 2.54$：出现雾气，直至
\omterm{ex:b1:phase-changes:dew}{相对湿度}降至 $2.08/3.17 = 66\%$ 以下。

\textbf{15.} $P = P_s(\qty{20}{\degreeCelsius}) = \qty{2.3}{kPa}$；$v = \qty{2.0e-3}{m^3/kg}$；
$v_v = RT/(MP_s) = 8.314 \times 293/(0.018 \times 2340) = \qty{58}{m^3/kg}$；$x =
(2.0 - 1.0) \times 10^{-3}/58 = 1.7 \times 10^{-5}$：蒸气 \qty{9}{mg}；
液体 $0.50 \times 10^{-3} = \qty{0.50}{L}$，占容器的一半。

\textbf{16.} $v = \qty{2.0e-3}{m^3/kg} < v_c$：等容线与饱和穹顶的液相支
相交——液体膨胀快于蒸发，液面
\emph{上升}，容器充满液体时 $v_l(T) =
\qty{2.0e-3}{m^3/kg}$，略高于 \qty{360}{\degreeCelsius}（约
\qty{190}{bar}）。此后，近似不可压缩的液体在定容下受热：压强
每\omterm{def:b1:kinetic-theory:temperature}{开尔文}升高数十巴——容器爆裂。
切勿密封充满液体的容器。

\textbf{17.} $v = \qty{1.0e-2}{m^3/kg} > v_c$：液面\emph{下降}；最后一滴
液体在 $v_v(T) = 0.010$ 时蒸发殆尽，温度介于 \qty{340}{\degreeCelsius}（$0.0108$）
与 \qty{350}{\degreeCelsius}（$0.0088$）之间：约 \qty{344}{\degreeCelsius}，
\qty{150}{bar}。

\textbf{18.} $v = \qty{3.2e-3}{m^3/kg} \approx v_c$：弯月面保持在
中部附近，随两相密度趋于相等，它越来越平、越来越淡，
流体变为乳白色（临界乳光），在 \qty{374}{\degreeCelsius}、
\qty{221}{bar} 时弯月面消失：只剩一相。

\textbf{19.} 三条竖直线：$v = 2.0 \times 10^{-3}$ 由
液相支穿出，位于 C 左侧；$10^{-2}$ 由蒸气支穿出，
位于 C 右侧；$3.2 \times 10^{-3}$ 穿过顶点。

\textbf{20.} 空气：\qty{0.5}{L}，\qty{1.0}{bar}，\qty{293}{K}；同体积下升至
\qty{373}{K}：$P_{\text{空气}} = 1.0 \times 373/293 = \qty{1.27}{bar}$；加上
$P_s = \qty{1.01}{bar}$：内部 \qty{2.3}{bar}，高于外界
\qty{1.3}{bar}，即 $1.3 \times 10^5 \times 10^{-2} = \qty{1.3}{kN}$ 作用于盖子上——相当于
\qty{130}{kg} 的重量。

\textbf{21.} \omterm{def:b1:phase-changes:metastable}{过冷}、亚稳：无\omterm{def:b1:phase-changes:metastable}{成核}位点（水洁净，
无尘埃、无扰动），故无法生成第一颗晶体。一次轻敲即可使冰\omterm{def:b1:phase-changes:metastable}{成核}，
冰在瓶中生长并释放 $L_f$；温度升至
\qty{0}{\degreeCelsius} 并停留于此，这是冰与
水在 \qty{1}{atm} 下共存的唯一温度：一旦
释放的热量使混合物达到此温度，凝固即告停止。

\textbf{22.} $x = c\,\Delta T/L_f = 4.18 \times 8/334 = 0.10$。

\textbf{23.} $\Delta s = c\ln(273/265) - xL_f/273 = 124.3 - 122.3 =
\qty{+2.0}{J/(kg.K)}$：\omterm{prop:b1:first-law:transformations}{绝热过程}，故全部为产生量；数值为正，符合
自发过程的要求。

\textbf{24.} $x = 4.18 \times 5/2257 = 0.93\%$：每升含水蒸气 \qty{9.3}{g}，
$9.3 \times 10^{-3} \times 1.67 = \qty{16}{L}$ 的蒸气在液体内部
于不足一秒的时间内生成：水杯喷涌而出。

\textbf{25.} $T_0S_{\text{产生}} = 273 \times 2.0 = \qty{550}{J}$，每升
与 $mgh = \qty{9.8}{J}$ 相比，是其五十六倍。
\omterm{def:b1:phase-changes:metastable}{亚稳态}储存了功——即可逆路径本可提取的那部分——而
一次轻敲便将其全部转化为熵。
\end{solution}
